% from ega2-1-fr.tex
\oldpage[II]{157}
donc définit un (TN)-isomorphisme de degré $0$ de $S$-modules gradués
$S_+(1)\to a\mathfrak{I}S$. Par ailleurs $x\mapsto a^{-1}x$ est un
isomorphisme de degré $0$ de $S$-modules gradués
$a\mathfrak{I}S\xrightarrow{\sim}\mathfrak{I}S$. On obtient donc ainsi par
composition (\hyperref[II.3.2.4-fr]{3.2.4}) un isomorphisme de
$\mathcal{O}_X$-Modules
$\mathcal{O}_X(1)\xrightarrow{\sim}\mathcal{I}\mathcal{O}_X$, et comme $S$
est engendrée par $S_1=\mathfrak{I}$, $\mathcal{O}_X(1)$ est inversible
(\hyperref[II.3.2.5-fr]{3.2.5}) et très ample
(\hyperref[II.4.4.3-fr]{4.4.3} et \hyperref[II.4.4.9-fr]{4.4.9}), d'où notre
assertion.

\subsection{Résultats préliminaires sur la localisation dans les anneaux gradués}
\label{subsection:II.8.2-fr}

\begin{env}[8.2.1]
\label{II.8.2.1-fr}
Soit $S$ un anneau gradué, où nous ne supposons pas pour le moment les degrés
positifs. On posera
\begin{equation}
\label{II.8.2.1.1-fr}
  S^{\geq}=\bigoplus_{n\geq0}S_n,
  \qquad
  S^{\leq}=\bigoplus_{n\leq0}S_n,
\tag{8.2.1.1}
\end{equation}
qui sont des sous-anneaux gradués de $S$, à degrés respectivement tous
positifs et tous négatifs. Si $f$ est un élément homogène de degré $d$
(positif ou négatif) de $S$, l'anneau des fractions $S_f=S'$ est encore muni
d'une structure d'anneau gradué, en prenant pour $S'_n$ ($n\in\mathbf{Z}$)
l'ensemble des $x/f^k$, où $x\in S_{n+kd}$ ($k\geq0$)~; on posera encore
$S_{(f)}=S'_0$, et on écrira $S_f^{\geq}$ et $S_f^{\leq}$ pour $S'^{\geq}$ et
$S'^{\leq}$ respectivement. Si $d>0$, on a
\begin{equation}
\label{II.8.2.1.2-fr}
  (S^{\geq})_f=S_f
\tag{8.2.1.2}
\end{equation}
puisque, si $x\in S_{n+kd}$, avec $n+kd<0$, on peut écrire
$x/f^k=xf^h/f^{h+k}$ et que $n+(h+k)d>0$ pour $h$ assez grand et $>0$. On en
conclut par définition que
\begin{equation}
\label{II.8.2.1.3-fr}
  (S^{\geq})_{(f)}=(S_f^{\geq})_0=S_{(f)}.
\tag{8.2.1.3}
\end{equation}

Si $M$ est un $S$-module gradué, on posera de même
\begin{equation}
\label{II.8.2.1.4-fr}
  M^{\geq}=\bigoplus_{n\geq0}M_n,
  \qquad
  M^{\leq}=\bigoplus_{n\leq0}M_n
\tag{8.2.1.4}
\end{equation}
qui sont respectivement un $S^{\geq}$-module gradué et un $S^{\leq}$-module
gradué et ont pour intersection le $S_0$-module $M_0$. Si $f\in S_d$, on
définit encore $M_f$ comme $S_f$-module gradué en prenant comme éléments de
degré $n$ les $z/f^k$, où $z\in M_{n+kd}$ ($k\geq0$)~; on désignera par
$M_{(f)}$ l'ensemble des éléments de degré $0$ de $M_f$, qui est un
$S_{(f)}$-module, et on écrira $M_f^{\geq}$ et $M_f^{\leq}$ au lieu de
$(M_f)^{\geq}$ et $(M_f)^{\leq}$ respectivement. Si $d>0$, on voit comme
ci-dessus que
\begin{equation}
\label{II.8.2.1.5-fr}
  (M^{\geq})_f=M_f
\tag{8.2.1.5}
\end{equation}
et
\begin{equation}
\label{II.8.2.1.6-fr}
  (M^{\geq})_{(f)}=(M_f^{\geq})_0=M_{(f)}.
\tag{8.2.1.6}
\end{equation}
\end{env}

\begin{env}[8.2.2]
\label{II.8.2.2-fr}
Soit $z$ une indéterminée, que nous nommerons \emph{variable
d'homogénéisation}. Si $S$ est un anneau gradué (à degrés positifs ou
négatifs), l'algèbre de polynômes\footnote{Il ne saurait ici y avoir de
confusion avec l'usage de la notation $\widehat{S}$ pour désigner le séparé
complété d'un anneau.}
\begin{equation}
\label{II.8.2.2.1-fr}
  \widehat{S}=S[z]
\tag{8.2.2.1}
\end{equation}

